%-------------------------
% Resume in Latex
% Author : Jake Gutierrez
% Based off of: https://github.com/sb2nov/resume
% License : MIT
%------------------------

\documentclass[letterpaper,10pt]{article}

\usepackage{latexsym}
\usepackage[empty]{fullpage}
\usepackage{titlesec}
\usepackage{marvosym}
\usepackage[usenames,dvipsnames]{color}
\usepackage{verbatim}
\usepackage{enumitem}
\usepackage[hidelinks]{hyperref}
\usepackage{fancyhdr}
\usepackage[english]{babel}
\usepackage{tabularx}
\usepackage{amssymb}
\usepackage{fontawesome5}
\usepackage{pifont}
\input{glyphtounicode}

\newcommand{\githubstars}[1]{[{\faGithubSquare}\,\ding{72}\,\textbf{#1}]}


%----------FONT OPTIONS----------
% sans-serif
% \usepackage[sfdefault]{FiraSans}
% \usepackage[sfdefault]{roboto}
% \usepackage[sfdefault]{noto-sans}
% \usepackage[default]{sourcesanspro}

% serif
% \usepackage{CormorantGaramond}
% \usepackage{charter}


\pagestyle{fancy}
\fancyhf{} % clear all header and footer fields
\fancyfoot{}
\renewcommand{\headrulewidth}{0pt}
\renewcommand{\footrulewidth}{0pt}

% Adjust margins
\addtolength{\oddsidemargin}{-0.55in}
\addtolength{\evensidemargin}{-0.55in}
\addtolength{\textwidth}{1.1in}
\addtolength{\topmargin}{-.65in}
\addtolength{\textheight}{1.3in}

\urlstyle{same}

\linespread{1.3}

\raggedbottom
\raggedright
\setlength{\tabcolsep}{0in}

% Sections formatting
\titleformat{\section}{
  \vspace{4pt}\scshape\raggedright\large
}{}{0em}{}[\color{black}\titlerule \vspace{3pt}]

% Ensure that generate pdf is machine readable/ATS parsable
\pdfgentounicode=1

%-------------------------
% Custom commands
\newcommand{\resumeItem}[1]{
  \item\small{
    {#1 \vspace{-2pt}}
  }
}

\newcommand{\resumeSubheading}[4]{
  \vspace{2pt}\item
    \begin{tabular*}{0.97\textwidth}[t]{l@{\extracolsep{\fill}}r}
      \textbf{#1} & #2 \\
      \textit{\small#3} & \textit{\small #4} \\
    \end{tabular*}\vspace{-3pt}
}

\newcommand{\resumeSubSubheading}[2]{
    \item
    \begin{tabular*}{0.97\textwidth}{l@{\extracolsep{\fill}}r}
      \textit{\small#1} & \textit{\small #2} \\
    \end{tabular*}\vspace{-3pt}
}

\newcommand{\resumeProjectHeading}[2]{
    \item
    \begin{tabular*}{0.97\textwidth}{l@{\extracolsep{\fill}}r}
      \small#1 & #2 \\
    \end{tabular*}\vspace{-1pt}
}

\newcommand{\resumeSubItem}[1]{\resumeItem{#1}\vspace{-1pt}}

\renewcommand\labelitemii{$\vcenter{\hbox{\tiny$\bullet$}}$}

\newcommand{\resumeSubHeadingListStart}{\begin{itemize}[leftmargin=0.15in, label={}]}
\newcommand{\resumeSubHeadingListEnd}{\end{itemize}}
\newcommand{\resumeItemListStart}{\begin{itemize}}
\newcommand{\resumeItemListEnd}{\end{itemize}\vspace{-0pt}}

%-------------------------------------------
%%%%%%  RESUME STARTS HERE  %%%%%%%%%%%%%%%%%%%%%%%%%%%%


\begin{document}

%----------HEADING----------
\begin{center}
    \textbf{\LARGE \scshape Zeyu Chen} \\ \vspace{1pt}
    \small
    \underline{zeyuc07@connect.hku.hk} $|$
    \href{https://zykev.github.io/}{\underline{zykev.github.io}} $|$
    \underline{+86 18260068713}
\end{center}


% Skills and Experience before Education is recommended for industry applications.
% Reference: https://www.sweresume.app/articles/faang-resume-guide/
%-----------EDUCATION-----------
\section{Education}
  \resumeSubHeadingListStart
    \resumeSubheading
      {The University of Hong Kong (HKU)}{Hong Kong}
      {Ph.D. in Artificial Intelligence}{Sep. 2024 -- present}
    \resumeSubheading
      {The University of Hong Kong (HKU)}{Hong Kong}
      {M.S. in Data Science}{Sep. 2022 -- Jan. 2024}
    \resumeSubheading
      {Nanjing University of Science and Technology (NJUST)}{Nanjing, China}
      {B.S. in Applied Statistics}{Sep. 2016 -- Jun. 2020}

  \resumeSubHeadingListEnd


%-----------EXPERIENCE-----------
% NOTE(hlim): Do NOT list Teaching Assistant (TA), grading, or academic service roles here.
% Industry hiring managers look for engineering and research impact in the Experience section.
% TA duties (office hours, grading, course prep) do not signal the skills tech companies hire for.
% To show mentorship or leadership instead, frame it as a bullet inside your research role:
%   e.g. "Mentored 5+ graduate researchers, contributing to 4+ published papers."
% \section{Experience}
%   \resumeSubHeadingListStart

%     \resumeSubheading
%       {Feed-forward 4D Clothed Human Reconstruction}{Feb. 2025 -- Present}
%       {Visual AI Lab, HKU (Advisor: Prof. Kai Han)}{Hong Kong}
%       \resumeItemListStart
%         \resumeItem{Leverage 3D Gaussian Splatting to reconstruct high-fidelity 3D human avatars from sparse-view inputs, optimizing for both geometric accuracy and rendering speed.}
%         \resumeItem{Develop a dual-layer decoupling pipeline to separate body and clothing representations, introducing a novel animation scheme via extended Linear Blend Skinning (LBS) to recover dynamic cloth deformations from monocular videos.}
%         \resumeItem{Achieve state-of-the-art (SOTA) performance, outperforming baselines such as \textit{LHM} and \textit{GHG} in terms of reconstruction quality, structural integrity, and visual fidelity.}
%       \resumeItemListEnd

%     \resumeSubheading
%       {Multimodal Unified Embedding Space Learning}{Jan. 2024 -- Sep. 2025}
%       {Visual AI Lab, HKU (Advisor: Dr. Jie Li, Prof. Kai Han)}{Hong Kong}
%       \resumeItemListStart
%         \resumeItem{Developed CodeBind, a unified framework to align nine distinct modalities (including image, text, audio, video, thermal, and depth) into a joint embedding space via Vector Quantization (VQ) techniques.}
%         \resumeItem{Proposed a decoupled representation learning scheme using modality-shared and modality-specific codebooks, effectively isolating cross-modal invariant features from domain-specific characteristics.}
%         \resumeItem{Surpassed industry-standard baselines (\textit{ImageBind, ViT-Lens}) across multiple benchmarks, achieving SOTA results in zero-shot multimodal classification and cross-modal retrieval tasks.}
%         \resumeItem{Conducted in-depth analysis of language density's role in multimodal alignment and leveraged Qwen3-VL to interpret and visualize the semantic properties of decomposed modality-specific components.}
%         \resumeItem{Demonstrated the versatility of the decoupled space by enabling downstream applications such as any-to-image generation.}
%       \resumeItemListEnd

%     \resumeSubheading
%       {Compositional Text-to-Image Generation with Diffusion Model}{Jun. 2024 -- Dec. 2024}
%       {Master Thesis, HKU (Advisor: Dr. Jie Li, Prof. Kai Han)}{Hong Kong}
%       \resumeItemListStart
%         \resumeItem{Developed a controllable generation framework based on Stable Diffusion to synthesize compositional images featuring multiple personalized concepts and precise layout arrangements.}
%         \resumeItem{Implemented LoRA-based fine-tuning integrated with an auxiliary image encoder to inject reference features into text embeddings, significantly enhancing identity preservation and training efficiency over the \textit{Mix-of-Show} baseline.}
%         \resumeItem{Engineered a novel training-free guidance mechanism by optimizing cross-attention maps during the sampling process, ensuring that the spatial distribution of personalized concepts strictly adheres to user-defined bounding boxes.}
%         \resumeItem{Demonstrated superior generalization in complex multi-concept scenarios, effectively mitigating the "attribute leakage" and "missing objects" issues common in standard T2I models.}
%       \resumeItemListEnd

    % \resumeSubheading
    %   {Machine Vision to Facial Fatigue with Nonlocal 3D Attention Network}{Feb. 2021 -- Sep. 2021}
    %   {Research Intern, BIAI Intelligence Biotech Co., Ltd (Advisor: Dr. Xiaotao Li)}{Shenzhen, China}
    %   \resumeItemListStart
    %     \resumeItem{Curated a real-world facial fatigue dataset by capturing high-resolution video data in naturalistic driving and working scenarios, moving beyond traditional lab-simulated environments.}
    %     \resumeItem{Developed a spatiotemporal feature extraction backbone by upgrading 2D architectures to a 3D ResNet integrated with Non-local Self-Attention mechanisms, achieving 90\% accuracy on binary fatigue classification.}
    %     \resumeItem{Enhanced model interpretability by utilizing Grad-CAM and guided backpropagation to visualize and validate salient physiological fatigue indicators (e.g., eye closure duration, blinking patterns) captured by the network.}
    %   \resumeItemListEnd

  % \resumeSubHeadingListEnd



%-----------RESEARCH INTERESTS-----------
\section{Research Interests}
  \small{Multimodal understanding, 3D vision, spatial intelligence, and medical image analysis.}


%-----------PUBLICATIONS-----------
\section{Publications}
 \begin{itemize}[leftmargin=0.15in, label={}]
    \item{
    \textbf{Chen, Z.}, Li, J., \& Han, K. (2026). CodeBind: Decoupled Representation Learning for Multimodal Alignment with Unified Compositional Codebook. In Findings of the Association for Computational Linguistics: ACL 2026. Association for Computational Linguistics.
    }
    \item{
    \textbf{Chen, Z.}, Liu, P., Han, K., Liao, P., Yang, Y., Wong, M. C. M., Yiu, C. K. Y., \& Lo, E. C. M. (2026). AI in Oral Health Surveillance: Critical Review. \textit{Journal of Dental Research (JDR)}.
    }
    % \item{
    % \textbf{Chen, Z.}, Zhang, X., Li, J., Ni, J., Chen, G., Wang, S., ... \& Li, X. (2021). Machine vision detection to daily facial fatigue with a nonlocal 3D attention network. arXiv preprint arXiv:2104.10420.
    % }
    % \item{
    % Ding, Z., \textbf{Chen, Z.}, Ma, M., \& Zhang, Z. (2019). Application of KF improved by Hermite interpolation in target trajectory correction. In AIP Conference Proceedings (Vol. 2185, No. 1, p. 020051). AIP Publishing LLC.
    % }
 \end{itemize}


%-----------PRESENTATIONS-----------
\section{Presentations}
  \begin{itemize}[leftmargin=0.15in, label={}]
    \small{\item{
      \textbf{DentMAG: Masked Autoencoding with Guidance Cues for Comprehensive Dental Diagnosis.} 2026 IADR/AADOCR/CADR General Session \& Exhibition, San Diego, CA, USA, March 23--28, 2026. \\
      \textbf{AI-assisted Photograph-based Oral Health Screening and Surveillance: A Scoping Review.} 2025 IADR/AADOCR/CADR General Session \& Exhibition, Barcelona, Spain, June 23--28, 2025.
    }}
  \end{itemize}


%-----------HONORS \& AWARDS-----------
\section{Honors \& Awards}
 \begin{itemize}[leftmargin=0.15in, label={}]
    \small{\item{
      \textbf{2024-2028} -- HKU Postgraduate Scholarship \\
     \textbf{2023} -- Suoxinda Scholarship in Data Science, HKU\\
    %  \textbf{2023} -- Subject Prize Awardees, HKU (HK\$1000 x4 courses) \\
    %  \textbf{2017-2019} -- NJUST Scholarship for Excellent Students
    }}
 \end{itemize}


%-----------TECHNICAL SKILLS-----------
\section{Technical Skills}
 \begin{itemize}[leftmargin=0.15in, label={}]
    \small{\item{
     \textbf{Programming}{: Python} \\
     \textbf{Frameworks \& Tools}{: PyTorch, OpenCV} \\
     \textbf{Languages}{: Chinese (Native), English (Fluent)}
    }}
 \end{itemize}


%-------------------------------------------
\end{document}
